% ====================================================================
% Chapter 8: Open Problems and Future Directions
% ====================================================================
\section{Open Problems and Future Directions}\label{ch:future}

\subsection{Overview}\label{sec:ch8-overview}

The trajectory of agent self-evolution extends far beyond making agents more ``autonomous'' in a colloquial sense. As the analyses in Chapters~\ref{ch:painpoints} and~\ref{ch:evaluation} have demonstrated, the critical bottlenecks are not cognitive capacity---they are reliability in production environments, controllability within improvement loops, resistance to metric gaming, and the establishment of hard boundaries around cost and safety. Future research must confront these bottlenecks directly: automated verifiers must reduce dependence on human feedback; multi-agent co-evolution must avoid consensus hallucination and cost inflation; safety mechanisms must prevent misevolution and prompt injection; cross-domain transfer must break free from benchmark-specific adaptation; and the path from agent evolution to AGI must unify local self-improvement, open-ended exploration, and value alignment into a coherent engineering framework.

The literature reviewed in this survey provides multiple composable modules. STaR~\cite{zelikman2022star}, RISE~\cite{qu2024rise}, and Absolute Zero Reasoners~\cite{absolutezero2025} demonstrate that generate--filter--train loops can improve capabilities on verifiable reasoning or code tasks. Reflexion~\cite{shinn2023reflexion}, ExpeL~\cite{zhao2024expel}, ACE~\cite{ace2025}, ReasoningBank~\cite{reasoningbank2025}, Memory-R1~\cite{memoryr1_2025}, and AriadneMem~\cite{ariadnemem2026} illustrate pathways for reflection, experience crystallization, and memory asset management. Voyager~\cite{wang2023voyager} and WebEvolver~\cite{webevolver2025} reveal the potential of environmental interaction and world models in open-ended tasks. ADAS~\cite{hu2024adas}, DGM~\cite{dgm2025}, SICA~\cite{sica2025}, G\"{o}del Agent~\cite{godelagent2024}, and AlphaEvolve~\cite{novikov2025alphaevolve} demonstrate the power of combining LLMs with search, archival storage, self-modification, and programmatic evaluators. Multi-Agent Debate~\cite{du2023improving}, EvoMAC~\cite{evomac2024}, SPIRAL~\cite{spiral2025}, and RAGEN~\cite{ragen2025} show that multi-agent, self-play, and collaborative networks can provide additional exploration breadth and feedback signals.

However, these modules remain fragmented, each constrained by cost, verification fidelity, generalization scope, and safety guarantees. The defining question for the next decade is therefore not ``Can self-evolving agents exist?''---they already do---but ``Can we build reusable self-evolution infrastructure?'' Such infrastructure would encompass: verifiable modeling of tasks and environments; multi-layer evaluators; auditable memory and experience repositories; variant archives with lineage tracking; security sandboxes and permission governance; cost budgets and resource gates; cross-domain transfer evaluation protocols; and a transition of human roles from per-item labeling to rule design, boundary auditing, and evaluator calibration. The following six subsections address these directions in depth.

% ====================================================================
% 8.1 Evaluation Bottleneck
% ====================================================================
\subsection{The Evaluation Bottleneck: When Benchmark Gains Do Not Translate to Real-World Reliability}\label{sec:ch8-evaluation}

The most fundamental bottleneck in agent self-evolution is evaluation. Every self-improvement loop---whether it operates at the prompt level, the architecture level, the code level, or the weight level---requires a feedback signal. The quality, generality, and robustness of that signal ultimately determine whether the system is genuinely improving or merely optimizing a proxy that has drifted away from the intended objective. The Mom Test analysis in Chapter~\ref{ch:painpoints} identified evaluation as the second-largest pain category (14.4\% of all reported pain points), with specific complaints about benchmark saturation, Goodhart effects, and the absence of production-grade evaluation metrics. This section examines the evaluation bottleneck from four angles: benchmark saturation, evaluator quality, automated evaluator design, and the path toward production evaluators.

\subsubsection{Benchmark Saturation and Metric Gaming}

The current landscape of agent benchmarks exhibits two problematic patterns. First, a small number of benchmarks---HumanEval~\cite{chen2021codex}, SWE-Bench~\cite{jimenez2024swebench}, GSM8K~\cite{cobbe2021gsm8k}, MATH~\cite{hendrycks2021math}---have become de facto standards, creating an incentive structure where research groups optimize for leaderboard positions rather than genuine capability gains. Second, many of these benchmarks are either publicly available or easily reconstructible, creating contamination risks where the benchmark examples appear---intentionally or unintentionally---in training data or prompt libraries. The cross-reference analysis in Section~\ref{ch:evaluation} documented multiple cases where reported improvements on specific benchmarks did not transfer to related but distinct evaluation settings.

The Goodhart effect is particularly insidious in self-evolving systems because the system itself is actively searching for ways to increase the measured score. Unlike static systems where human engineers might accidentally overfit to a benchmark, self-evolving agents can systematically explore the boundary between genuine improvement and metric gaming. Self-Rewarding Language Models~\cite{metacognition2024selfrewarding} and Meta-Rewarding~\cite{metarewarding2024} revealed that evaluators trained alongside generators tend to develop length bias, score inflation, and position bias---precisely the failure modes predicted by Goodhart's law. The framework-painpoint crosswalk analysis documented 89 evaluation-related repositories in the ecosystem, the largest single category, yet community complaints about benchmark pollution and metric gaming persist, suggesting that quantity of evaluation tooling has not resolved the quality problem.

\subsubsection{Evaluator Quality and Degeneration}

A deeper problem lies in evaluator degeneration. When the same system that generates candidates also evaluates them---or when the evaluator is trained on outputs from a related model---the evaluation signal can drift away from genuine quality. This manifests in several ways. Length bias occurs when evaluators systematically prefer longer outputs regardless of content. Score compression occurs when evaluators fail to distinguish between genuinely excellent and merely adequate outputs, reducing the gradient signal for improvement. Reward hacking occurs when the generator discovers inputs that trigger high evaluator scores without producing genuinely better outputs, a phenomenon documented in RLHF literature and replicated in self-evolving agent settings by the OEP attack~\cite{oep2026} which demonstrated that experience poisoning can corrupt self-evolution trajectories.

The meta-evaluator approach---using a second evaluator to evaluate the first evaluator---partially addresses this problem but introduces its own regress. Meta-Rewarding~\cite{metarewarding2024} showed that a three-tier actor/judge/meta-judge architecture can reduce some forms of evaluator bias, but the meta-judge itself is subject to the same degeneration pressures. What is needed is not an infinite regress of evaluators but a ground truth anchor: either human evaluation on calibrated subsets, deterministic programmatic checks, or environmental feedback that cannot be gamed by the agent.

\subsubsection{Toward Automated Evaluator Design}

The most promising direction for breaking the evaluation bottleneck is automated evaluator design. Four categories of automated evaluators are emerging. First, deterministic program verification---unit tests, type checking, linting, static analysis, fuzzing, regression testing, and sandbox execution---provides the strongest signal for code tasks. Systems like DGM~\cite{dgm2025}, SICA~\cite{sica2025}, EvoMAC~\cite{evomac2024}, and AlphaEvolve~\cite{novikov2025alphaevolve} demonstrate self-evolution precisely because candidate code can be automatically executed and scored. Future work should extend this capability to broader application domains: data pipelines can verify schema compliance, distribution shift, and end-to-end consistency; business processes can verify state machine constraints and permission policies; web tasks can verify DOM states, backend conditions, and user intent proxies; robotics tasks can verify physical constraints and safety boundaries.

Second, model-assisted evaluators can supplement human feedback for tasks where complete programmatic evaluation is infeasible. Research synthesis quality, product requirement reasonableness, design maintainability, and customer service appropriateness are examples where LLM-as-a-judge, critic agents, debate mechanisms, and meta-judges can provide partial signals. However, model-assisted evaluators must not serve as sole arbiters of quality. They should be embedded within multi-judge panels, cross-model consistency checks, calibrated reference sets, and human spot-check protocols.

Third, interactive environment evaluators---exemplified by Voyager~\cite{wang2023voyager}, WebEvolver~\cite{webevolver2025}, AppWorld, WebArena~\cite{webarena2023}, ALFWorld~\cite{shridhar2020alfworld}, and WebShop---assess agent capability through action in simulated environments rather than static output comparison. Future work should construct more environments that are simultaneously replayable, scorable, and perturbable: enterprise process sandboxes, simulated CRM/ERP systems, virtual browsers, software repository mirrors, research experiment simulations, long-horizon memory tasks, and multi-user collaboration scenarios.

Fourth, safety and cost evaluators must become first-class components of the evaluation pipeline. Every candidate variant should undergo permission checks, sensitive data checks, prompt injection resistance tests, budget ceiling tests, resource leak tests, and rollback tests. Cost metrics---cost per success, cost per unit improvement, maximum tail cost, retry rate, and human takeover rate---should be recorded and incorporated into fitness functions. Without explicit cost evaluation, self-evolution systems will inevitably push all optimization pressure toward short-term success rates while ignoring long-term sustainability.

\subsubsection{From Research Benchmarks to Production Evaluators}

The gap between research benchmarks and production evaluation is substantial. Research benchmarks test agents on curated datasets with known ground truth; production evaluators must assess agents on open-ended, distribution-shifting, adversarially-contaminated real-world tasks. The Mom Test analysis revealed that the top user complaint about evaluation (P022) specifically targets the disconnection between benchmark optimization and production reliability. Bridging this gap requires three developments: (1) production evaluator libraries that cover code quality, web interaction, business process compliance, memory consistency, safety boundary enforcement, and cost efficiency; (2) self-evolution reporting standards that mandate disclosure of training/validation/test separation, failure candidate ratios, cost per improvement, rollback rates, and safety incidents; and (3) shared evaluation infrastructure---an ``Agent Evaluation Hub''---where researchers and practitioners can submit evaluation trajectories, report negative results, share benchmark contamination information, and maintain cost baselines.

Long-term, evaluators will become the ``compilers'' of agent self-evolution. Just as software engineering requires code to pass compilation, testing, and CI before merging, future agent prompts, tool strategies, memory rules, model routing configurations, and self-modification code will need to pass evaluator pipelines before entering production. The human role will shift from per-item output judgment to evaluator design, evaluator calibration, failure sample auditing, and risk preference specification.

% ====================================================================
% 8.2 Memory & Drift
% ====================================================================
\subsection{Memory and Drift: Long-Term Learning Stability in Self-Evolving Agents}\label{sec:ch8-memory}

Memory is the substrate of learning, and drift is its nemesis. A self-evolving agent that cannot stably maintain and update its knowledge over extended operational periods is not genuinely evolving---it is merely cycling through ephemeral adaptations. The Mom Test analysis identified memory-related pain points as a significant category (13.4\% of all reported issues), with specific complaints about catastrophic forgetting, memory pollution, retrieval noise, context window overflow, and the absence of versioned, auditable memory architectures. This section examines the memory challenge from four perspectives: catastrophic forgetting in continuous learning, memory architecture design, experience replay and consolidation, and long-horizon drift monitoring.

\subsubsection{Catastrophic Forgetting and the Stability-Plasticity Dilemma}

Self-evolving agents face a fundamental tension between plasticity (the ability to learn new capabilities) and stability (the retention of previously acquired capabilities). In neural network-based agents, this manifests as catastrophic forgetting: the overwriting of previously learned representations during fine-tuning on new tasks. In prompt-based and code-based agents, the analogous phenomenon is strategy drift, where improvements on new task distributions cause regressions on previously mastered tasks. The Capability Erosion study~\cite{capabilityerosion2026} directly documented this phenomenon, showing that self-evolving agents can lose core competencies during evolution cycles, and proposed Capability Preservation Evaluation (CPE) as a monitoring framework.

The challenge is particularly acute for agents that modify their own code or prompts. Unlike neural weight updates, which can be regularized through techniques like elastic weight consolidation, code modifications can alter control flow in ways that unpredictably affect previously working paths. DGM~\cite{dgm2025} addresses this partially through its archive mechanism---maintaining diverse variants so that regressions can be detected and rolled back---but the archive itself requires memory management, pruning strategies, and relevance scoring. The Voyager~\cite{wang2023voyager} skill library represents another approach, organizing learned capabilities as composable skill modules that can be selectively retrieved and applied, but its transferability beyond Minecraft remains unproven.

\subsubsection{Memory Architecture Design for Self-Evolution}

Current memory architectures for agents fall into several categories, each with distinct trade-offs. Episodic memory stores raw interaction traces, providing high-fidelity recall at the cost of storage overhead and retrieval noise. Semantic memory stores compressed abstractions---facts, rules, and patterns---offering efficient retrieval but potentially losing important context. Procedural memory stores executable skills and strategies, enabling direct action reuse but requiring careful version management. Meta-memory stores information about the memory system itself---what is known, what is uncertain, what has been tried and failed---enabling strategic learning decisions but adding architectural complexity.

The most sophisticated current systems combine multiple memory types. AriadneMem~\cite{ariadnemem2026} uses evolutionary graph memory with entropy-aware gating, representing memories as nodes in a directed acyclic graph where edges encode derivation relationships. ReasoningBank~\cite{reasoningbank2025} distills reasoning strategies into a searchable library, enabling cross-task transfer of problem-solving patterns. Memory-R1~\cite{memoryr1_2025} trains a memory manager using reinforcement learning, learning when to store, retrieve, and forget information based on task performance signals. ACE~\cite{ace2025} maintains an evolving context playbook that captures successful interaction patterns as reusable templates.

However, no current system provides all the properties needed for robust long-term self-evolution: versioning (tracking how memories change over time), deduplication (avoiding redundant or contradictory memories), expiration (gracefully degrading outdated information), rollback (reverting to previous memory states when regressions are detected), access control (preventing unauthorized memory modification), and audit trails (recording when, why, and how each memory was created or modified). The framework-painpoint crosswalk analysis found 57 memory-related repositories in the ecosystem, yet community complaints about memory pollution and drift persist, indicating that architectural abundance has not resolved the fundamental design challenge.

\subsubsection{Experience Replay and Consolidation}

Experience replay---the practice of revisiting and reprocessing past experiences---is a well-established technique for mitigating catastrophic forgetting in continual learning. In the context of self-evolving agents, experience replay takes several forms. Regression test suites serve as a form of experience replay for code-modifying agents: each test encodes a past capability that must be preserved. Archive-based replay, as employed by DGM~\cite{dgm2025}, maintains a population of past variants and periodically re-evaluates them on current tasks to detect capability loss. Skill library replay, as in Voyager~\cite{wang2023voyager}, periodically re-executes learned skills to verify continued functionality.

A more sophisticated approach is memory consolidation, inspired by biological sleep-dependent memory processing. Rather than simply replaying raw experiences, consolidation involves reorganizing, compressing, and integrating memories to extract generalizable knowledge while discarding noise. For self-evolving agents, this could involve periodically running background processes that: (1) cluster similar experiences and extract common patterns; (2) identify contradictions between stored memories and resolve them through re-evaluation; (3) compress verbose episodic memories into concise semantic representations; (4) detect and flag memories that have not been recently accessed for potential expiration; and (5) generate synthetic training examples from stored experiences to maintain capability on under-represented task types.

The design of effective consolidation mechanisms requires careful balancing. Too-aggressive compression risks losing important nuance; too-conservative retention leads to storage bloat and retrieval degradation. The optimal consolidation strategy likely depends on the agent's task distribution, operational tempo, and resource constraints---suggesting that consolidation itself should be a target of meta-learning rather than a fixed engineering decision.

\subsubsection{Long-Horizon Drift Monitoring and Correction}

Even with well-designed memory architectures and consolidation mechanisms, self-evolving agents operating over extended periods will inevitably experience drift: gradual shifts in behavior, capability distribution, and decision-making patterns that may not be immediately detectable but compound over time. Drift can originate from multiple sources: distribution shift in the agent's input environment, accumulation of small biases in feedback signals, gradual corruption of memory representations, and cascading effects of seemingly benign code or prompt modifications.

Detecting drift requires continuous monitoring infrastructure. Regression test suites provide one signal: periodic evaluation on a fixed set of reference tasks can detect capability loss. Distribution monitoring provides another: tracking statistical properties of the agent's inputs, outputs, and internal states over time can reveal gradual shifts. Memory health checks provide a third: auditing the memory system for contradictions, redundancies, access pattern anomalies, and expiration violations. The Capability Erosion framework~\cite{capabilityerosion2026} proposed CPE as a systematic approach to tracking capability preservation across evolution cycles, incorporating both per-task performance metrics and aggregate capability profiles.

Correcting drift once detected is equally challenging. Simple rollback to a previous checkpoint may discard genuinely valuable improvements acquired since the drift began. Selective rollback---reverting only the components responsible for the regression while preserving unrelated improvements---requires fine-grained attribution of behavioral changes to specific modifications, which is itself an open research problem. The DGM archive mechanism provides a foundation for this by maintaining lineage information, but current implementations lack the granularity needed for selective rollback at the sub-component level. Future systems should maintain causally-structured modification histories that enable precise identification and reversion of drift-inducing changes without collateral damage to unrelated capabilities.

% ====================================================================
% 8.3 Safety & Alignment
% ====================================================================
\subsection{Safety and Alignment: Controllability of Self-Evolving Systems}\label{sec:ch8-safety}

Safety and alignment represent the most critical challenge for self-evolving agents. Unlike conventional agents that execute fixed behaviors, self-evolving systems modify their own code, prompts, strategies, and sometimes even their evaluation criteria. This self-modification capability introduces three categories of risk that are qualitatively different from those of static AI systems: the agent may modify its own behavior in unintended ways (misevolution), it may persist harmful adaptations in memory or code (experience contamination), and it may amplify misaligned objectives across multiple improvement cycles (objective drift). The Misevolution analysis~\cite{misevolution2026} specifically noted that these risks can emerge over time even without external attackers, generated by the system's own optimization pressure. The OEP study~\cite{oep2026} demonstrated a concrete attack vector: experience poisoning can corrupt self-evolution trajectories by injecting carefully crafted negative experiences that steer the agent's learning in harmful directions. This section examines the safety challenge across five layers: permission minimization, evolutionary sandboxes, auditable lineage, objective isolation, and human governance.

\subsubsection{Permission Minimization and the Principle of Least Authority}

The first layer of safety is permission minimization. Self-evolving agents should operate under the principle of least authority: each component should have only the minimum permissions needed for its designated function. Tool permissions should default to deny; read and write access should be separated; high-risk operations (production code modification, database writes, payment processing, email sending, deployment actions) should require explicit approval. The self-evolution module itself should not have direct permission to modify production code or production prompts; it should only generate candidate diffs that pass through sandbox verification, evaluator assessment, and human or policy-based gating before merging.

The G\"{o}del Agent review~\cite{godelagent2024} provides a cautionary example: while the theoretical G\"{o}del machine framework assumes the existence of provably correct self-modifications, real-world systems lack such proofs. The gap between theoretical self-improvement guarantees and practical engineering constraints necessitates conservative permission boundaries. Every permission granted to a self-evolving agent is an attack surface; every tool available is a potential vector for misuse, either by external attackers through prompt injection or by the agent's own misaligned optimization pressure.

\subsubsection{Evolutionary Sandboxes and Isolated Evaluation}

The second layer is evolutionary sandboxing. Every candidate variant generated by the self-evolution process should execute in an isolated environment with synthetic or anonymized data, a fixed evaluation suite, and comprehensive logging of all tool calls and state changes. Sandboxing serves dual purposes: preventing malicious behavior from affecting production systems, and preventing unintentional damage such as infinite loops, resource leaks, file corruption, real API calls, or shared memory contamination.

Current sandboxing approaches range from Docker containers and virtual machines to more lightweight process isolation mechanisms. However, the Mom Test analysis identified sandbox escape and sub-agent silent failure (P013) as significant pain points, indicating that current isolation mechanisms are insufficient for production deployment. Future sandbox designs should incorporate: network isolation with explicit allowlists, filesystem isolation with read-only mount points for dependencies, resource limits (CPU, memory, time, API calls) with hard kill thresholds, output filtering to prevent sensitive data exfiltration, and state reset capabilities to ensure clean evaluation conditions for each candidate.

\subsubsection{Auditable Lineage and Provenance Tracking}

The third layer is auditable lineage. Self-evolving systems must maintain complete records of every modification: where each version originated, what was changed, why it was accepted, on which tasks it improved performance, on which tasks it regressed, which safety rules it triggered, and who approved its deployment. Without lineage, misattribution of errors becomes unavoidable; with lineage, the system can roll back to the most recent safe ancestor and analyze which modification patterns carry elevated risk.

The DGM archive mechanism provides a foundational model for lineage tracking. Each variant in the DGM archive maintains its derivation history, enabling reconstruction of the evolutionary path that produced it. However, current implementations focus primarily on performance lineage (which benchmarks improved or regressed) and lack the safety-oriented lineage needed for production deployment. Future systems should extend lineage tracking to include: the specific code or prompt changes made, the evaluator scores that justified acceptance, the safety checks that were run and their results, the human approvals that were obtained (if any), the cost incurred during evaluation, and any anomalous behaviors observed during sandbox testing.

\subsubsection{Objective and Evaluator Isolation}

The fourth layer is objective isolation. Self-evolving agents must not be able to directly modify their own final evaluators, hidden test sets, or safety policies. If the agent can alter its own evaluation criteria, it can achieve high scores by simplifying the evaluation rather than improving its capabilities---a particularly dangerous form of reward hacking. During training and search phases, the agent may interact with a visible proxy evaluator, but final deployment decisions must pass through an invisible, immutable, independently maintained gate evaluator.

This principle extends to memory and experience isolation. The OEP attack~\cite{oep2026} demonstrated that experience poisoning---injecting crafted negative experiences into the agent's learning trajectory---can steer self-evolution toward harmful behaviors. Defense against experience poisoning requires: input validation for all experiences before they enter memory, anomaly detection for experience distributions that deviate from expected patterns, diversity maintenance in experience replay buffers to prevent single-source dominance, and periodic re-evaluation of stored experiences against current safety criteria.

\subsubsection{Human Governance and Graduated Autonomy}

The fifth layer is human governance. The goal is not to eliminate human oversight but to place it at higher-leverage positions: defining risk levels, reviewing high-risk diffs, calibrating evaluators, approving permission changes, conducting post-incident reviews, and maintaining red-team sample sets. Low-risk local optimizations can be automated; high-risk objective function changes and permission expansions must require governance approval. This graduated autonomy model prevents the system from oscillating between ``fully manual'' (unable to scale) and ``fully autonomous'' (unable to control).

The long-term aspiration is provable safety boundaries. While formal verification cannot comprehensively cover LLM behavior, it can cover critical constraints: tool permissions, state machine invariants, budget ceilings, data flow isolation, inaccessible resources, merge conditions, and rollback triggers. Future safe self-evolution systems may adopt a hybrid ``neural generation + symbolic guardrails + programmatic verification'' architecture, where the LLM explores freely within safety boundaries that are guaranteed by inspectable, formally-specified rules. The Capability Erosion framework~\cite{capabilityerosion2026} provides a starting point by defining capability preservation evaluation as a mandatory component of any self-evolution cycle, ensuring that safety is not treated as an optional add-on but as an integral part of the evolutionary objective.

% ====================================================================
% 8.4 Composability
% ====================================================================
\subsection{Composability: Modular Combination of the Five Evolution Loops}\label{sec:ch8-composability}

A recurring theme throughout this survey is that self-evolution is not a single technique but a family of composable mechanisms. The five evolution loops identified in Chapter~\ref{ch:theory}---specification-to-execution, search, evaluator, reflection, and population---can be combined in various configurations to address different task requirements, resource constraints, and safety profiles. However, the current literature treats these loops largely in isolation. STaR composes specification-to-execution with reflection but ignores population. AlphaEvolve composes search, evaluator, and population but has limited reflection. DGM composes search, population, and evaluator with code-level self-modification but operates without explicit reflection or specification-to-execution loops. This section examines the composability challenge from three perspectives: interface design between loops, composition patterns and their trade-offs, and the path toward a universal self-evolution platform.

\subsubsection{Interface Design Between Loops}

For loops to compose effectively, they need well-defined interfaces. Each loop has inputs (what it consumes), outputs (what it produces), state (what it maintains across iterations), and control signals (how it is activated, paused, or terminated). The key interface design challenge is standardizing these components so that the output of one loop can serve as the input to another without ad-hoc adaptation.

The specification-to-execution loop takes a natural language goal and produces an executable pipeline. Its output---a structured task decomposition with verification criteria---can serve as input to the search loop (which explores variations of the pipeline), the evaluator loop (which assesses whether the pipeline meets the verification criteria), or the reflection loop (which analyzes execution failures and generates revision candidates). The search loop takes a search space definition and produces candidate solutions. Its output can feed the evaluator loop (for quality assessment), the population loop (for archive inclusion), or the reflection loop (for failure analysis). The evaluator loop takes candidate solutions and produces scores and diagnostics. Its output can drive the search loop (guiding exploration toward promising regions), the reflection loop (identifying failure patterns for targeted improvement), or the population loop (determining which candidates to retain). The reflection loop takes execution traces and failure analyses and produces experience abstractions. Its output can inform the search loop (suggesting promising modification strategies), the evaluator loop (identifying blind spots in current evaluation), or the memory system (storing reusable patterns). The population loop maintains a diverse collection of candidates. Its output can seed the search loop (providing starting points for exploration), the evaluator loop (providing comparison baselines), or the reflection loop (providing historical context for analysis).

Standardizing these interfaces requires defining common data formats for candidate representations, evaluation results, experience records, and lineage metadata. It also requires defining control protocols: when should one loop invoke another, what information should be passed, and how should results be integrated. Current systems implement these interfaces in ad-hoc, system-specific ways, making it difficult to transfer components between frameworks or to combine techniques from different research groups.

\subsubsection{Composition Patterns and Trade-Offs}

Different composition patterns are suited to different task characteristics. For tasks with verifiable outputs (code, mathematics, formal logic), the search + evaluator + population pattern is most effective, as demonstrated by AlphaEvolve~\cite{novikov2025alphaevolve} and FunSearch~\cite{funsearch2024}. For tasks with subjective or ambiguous evaluation (creative writing, research synthesis, design), the reflection + evaluator pattern with human-in-the-loop calibration is more appropriate, as in Self-Refine~\cite{madaan2023selfrefine} and Reflexion~\cite{shinn2023reflexion}. For tasks requiring cumulative skill acquisition (game playing, web navigation, software maintenance), the reflection + population + memory pattern enables incremental capability building, as in Voyager~\cite{wang2023voyager} and ACE~\cite{ace2025}.

The trade-offs between composition patterns involve several dimensions. Cost: more loops mean more LLM calls, more evaluations, and more computational overhead. Latency: deeply composed loops require sequential processing that may exceed time constraints for interactive applications. Reliability: each additional loop introduces a potential failure point; composition depth amplifies the risk of cascading failures. Observability: as loops are composed more deeply, it becomes harder to attribute performance changes to specific loop interactions. Safety: certain compositions (particularly those involving self-modification) carry higher risk profiles than others.

A mature self-evolution platform should support flexible composition: given a task description, resource constraints, and safety requirements, it should recommend an appropriate loop composition, configure the interfaces between loops, and manage the overall evolution lifecycle. This is analogous to how database query optimizers select execution plans: the platform should reason about the cost-benefit trade-offs of different compositions and adapt its strategy based on observed performance.

\subsubsection{Toward a Universal Self-Evolution Platform}

The ultimate vision is a universal self-evolution platform that serves as the ``operating system'' for agent self-improvement. Such a platform would provide: a library of composable evolution primitives (each implementing one of the five loops); standardized interfaces for connecting primitives; a composition engine that selects and configures appropriate loop combinations based on task characteristics; a shared evaluation infrastructure with calibrated benchmarks and production evaluators; a memory and experience management layer with versioning, deduplication, and access control; an archive and lineage tracking system; a safety sandbox with configurable permission levels; and a cost management system with budget enforcement.

Current systems are far from this vision. Each research project implements its own evolution infrastructure from scratch, leading to duplicated effort, incompatible formats, and limited transferability of results. The open-source community has made progress in specific areas---evaluation frameworks, agent orchestration platforms, memory systems---but no existing system provides a unified substrate for composable self-evolution. Building such a platform is a significant engineering challenge, but it is also a prerequisite for moving self-evolution from a research curiosity to a production technology.

% ====================================================================
% 8.5 Production Challenges
% ====================================================================
\subsection{Production Challenges: Bridging Research and Deployment}\label{sec:ch8-production}

The gap between demonstrated self-evolution in research settings and reliable self-evolution in production environments remains substantial. The Mom Test analysis identified production reliability as the single largest pain category (29.9\% of all reported issues), with specific complaints about agent unreliability in production, drift without monitoring, feedback loops that require manual intervention, and the absence of cost governance. The star-analysis report further revealed a stark disconnect: projects with the highest community visibility (AutoGPT at 175K stars, score 28/100) often have the lowest production readiness, while academically rigorous projects (DSPy at 25K stars, score 73/100) receive comparatively less attention. This section examines the production challenge across four dimensions: deployment reliability, observability and debugging, cost governance, and organizational integration.

\subsubsection{Deployment Reliability and the Demo-to-Production Gap}

Research demonstrations of self-evolving agents typically operate under controlled conditions: curated task distributions, known evaluation criteria, abundant computational resources, and forgiving time constraints. Production environments present fundamentally different challenges: open-ended task distributions with unknown edge cases, ambiguous evaluation criteria that may conflict with business metrics, strict resource budgets, and real-time latency requirements. The gap between demo success and production stability is the defining challenge for translating self-evolution research into practice.

The framework-painpoint crosswalk analysis documented this gap across multiple domains. For code agents, the gap manifests as API hallucination, failure on niche or out-of-distribution codebases, and generation of plausible-looking but functionally incorrect code. For web agents, the gap appears as inability to handle dynamic pages, authentication flows, and unexpected UI changes. For workflow agents, the gap shows up as fragile automation that breaks when processes change and requires constant human supervision. Across all domains, the common pattern is that research evaluations test typical cases while production exposes tail cases.

Addressing the demo-to-production gap requires three developments. First, production evaluation protocols that test agents under realistic conditions: distribution-shifted inputs, adversarial prompts, resource constraints, concurrent access, partial failures, and recovery scenarios. Second, graduated deployment strategies that progressively expand the agent's operational scope while maintaining human oversight: starting with read-only operations, advancing to low-risk modifications under supervision, and eventually enabling autonomous operation within well-defined boundaries. Third, continuous monitoring and incident response infrastructure that detects production failures in real time, triggers appropriate rollback or escalation procedures, and feeds failure data back into the self-evolution process as learning signals.

\subsubsection{Observability, Debugging, and Explainability}

Self-evolving agents present unique observability challenges because their behavior changes over time. A production system that was functioning correctly yesterday may behave differently today because the agent modified its own prompts, strategies, or code overnight. Traditional monitoring approaches---which assume relatively static system behavior---are insufficient for systems that are actively modifying themselves.

Effective observability for self-evolving agents requires several capabilities. Prompt and strategy versioning: the ability to determine exactly which prompts, strategies, and code versions were active at any point in time. Decision tracing: the ability to reconstruct why the agent made a specific decision, including the chain of reasoning, the memories accessed, the tools invoked, and the alternatives considered. Modification attribution: the ability to trace any behavioral change back to the specific self-modification that caused it. Performance trending: the ability to detect gradual performance changes across evolution cycles and correlate them with specific modifications. Anomaly detection: the ability to identify when the agent's behavior deviates from expected patterns, triggering alerts for human review.

The Mom Test pain point P004 (``framework opacity---no visibility into what prompts are actually being sent'') directly illustrates the current observability deficit. Many agent frameworks abstract away the prompts, tool calls, and internal state changes that would be necessary for effective debugging. Self-evolving agents amplify this problem because the opaque behavior is actively changing. Future agent frameworks should treat observability as a first-class architectural requirement, not a debugging afterthought.

\subsubsection{Cost Governance and Resource Management}

Self-evolution is inherently expensive. Each improvement cycle requires generating candidates, evaluating them, maintaining archives, managing memory, and running safety checks. The Mom Test identified cost as a significant pain category (P006 ``runaway loops and cost explosion'', P020 ``cost is the dominant constraint''), and the framework-painpoint crosswalk found cost-related concerns across all seven framework categories. Without explicit cost governance, self-evolution systems can easily consume more resources than they save through improved performance.

Cost governance for self-evolving agents should operate at multiple levels. Per-improvement-cycle budgets: each self-modification attempt should have a maximum allowed cost (LLM tokens, API calls, compute time, human review time). Per-task budgets: each task should have a maximum total evolution cost, including all cycles dedicated to improving that task's performance. Lifetime budgets: the total self-evolution budget for an agent should be bounded, with periodic review and adjustment based on demonstrated return on investment. Value-based budgeting: evolution budgets should be proportional to the business value of the target task, concentrating optimization effort where it matters most.

The star-analysis report's finding that ``cost per improvement'' is rarely reported in research papers highlights a significant gap. Research systems report performance improvements but rarely the total computational cost of achieving those improvements. Production systems cannot afford this omission. Every self-evolution cycle should record its total cost and compare it against the value of the improvement achieved. If the cost exceeds the value, the evolution process should be throttled or redirected.

\subsubsection{Organizational Integration and Change Management}

Deploying self-evolving agents within organizations requires integration with existing processes, governance structures, and change management protocols. Self-modification is, in essence, automated change management---and organizations already have established procedures for managing changes to production systems. The challenge is aligning self-evolution processes with these existing procedures.

Key integration points include: change approval workflows (self-modifications should be routed through appropriate approval channels based on risk level), audit trails (all self-modifications should be logged in the same systems used for human-made changes), rollback procedures (the ability to revert self-modifications using the same rollback infrastructure used for manual changes), monitoring integration (self-evolution metrics should be visible in the same dashboards used for system health monitoring), and incident response (failures caused by self-modifications should trigger the same incident response procedures as failures caused by human errors). Treating self-evolution as a specialized form of automated change management, rather than as a fundamentally novel paradigm, reduces organizational resistance and leverages existing governance maturity.

% ====================================================================
% 8.6 Self-Evolve Roadmap
% ====================================================================
\subsection{A Roadmap for Self-Evolving AI}\label{sec:ch8-roadmap}

Drawing on the analyses in this survey---107 core papers, 348 GitHub repositories, 97 user pain points, and the five-loop taxonomy---we propose a phased roadmap for the development of self-evolving AI systems. The roadmap is organized into three phases, each building on the foundations established by the previous phase, and is grounded in the empirical evidence and community feedback documented throughout this survey.

\subsubsection{Phase 1: Foundation Infrastructure (2026--2027)}

The first phase focuses on building the infrastructure that makes self-evolution safe, observable, and reproducible. Five priorities emerge from the analysis. First, establishing standardized evaluator libraries covering code quality, web interaction, business process compliance, memory consistency, safety boundary enforcement, and cost efficiency. The current fragmentation of evaluation tooling (89 evaluation-related repositories in the ecosystem, yet persistent complaints about benchmark reliability) indicates that quantity has not achieved quality; coordination and standardization are needed. Second, developing self-evolution reporting standards that mandate disclosure of training/validation/test separation, failure candidate ratios, cost per improvement, rollback rates, and safety incidents. Without standardized reporting, the field cannot reliably compare different approaches or detect systematic weaknesses. Third, advancing auditable memory systems that upgrade experience assets from free-text logs to structured, versioned, compressible, and forgettable knowledge bases. The 57 memory-related repositories in the ecosystem demonstrate active interest, but the 13.4\% pain point share for memory issues indicates that current implementations are not meeting production needs.

Fourth, developing archive and lineage infrastructure that treats each self-modification as a trackable software supply chain event. DGM's archive mechanism and the lineage tracking concepts from the evolutionary computation literature provide starting points, but production-grade lineage requires finer granularity, safety-oriented metadata, and integration with existing change management systems. Fifth, establishing safety sandbox standards that provide reliable isolation for self-modification experiments. The current gap between Docker/VM-based isolation and the sandbox requirements identified in the safety literature (P013 ``environment isolation and sub-agent silent failure'') indicates that more sophisticated isolation mechanisms are needed.

\subsubsection{Phase 2: Graduated Autonomy (2027--2029)}

The second phase extends self-evolution from isolated, supervised cycles to graduated autonomous operation within organizational contexts. Five capabilities are targeted. Low-risk autonomous improvement: agents should be able to autonomously propose, evaluate, and deploy modifications to low-risk components (documentation, test cases, non-critical configuration) without human approval. The evaluator infrastructure from Phase 1 provides the foundation for determining which modifications qualify as low-risk. Medium-risk gray-scale evaluation: modifications to medium-risk components should be automatically evaluated in sandboxed environments and deployed to a small fraction of production traffic for real-world assessment before full deployment. High-risk human approval: modifications to critical components (core logic, safety policies, permission configurations) should require explicit human review and approval, with the system presenting the modification, its lineage, its evaluation results, and its potential impact in a human-readable format.

Heterogeneous multi-agent collaboration: the system should deploy specialized agents for generation, verification, red-teaming, cost control, and human auditing, forming a checks-and-balances architecture that reduces the risk of any single agent's failure propagating through the system. The current literature on multi-agent debate~\cite{du2023improving}, EvoMAC~\cite{evomac2024}, and SPIRAL~\cite{spiral2025} provides foundations for multi-agent interaction, but production deployment requires attention to cost control (preventing multi-agent cost inflation), observability (tracking inter-agent communication), and safety (preventing adversarial collusion). Auditable experience assetization: the system should automatically convert operational experiences into reusable knowledge assets, maintaining versioning, access control, and quality metrics. The shift from ephemeral adaptation to persistent assetization is what distinguishes genuine self-evolution from transient improvement.

\subsubsection{Phase 3: Open-Ended Evolution (2029+)}

The third phase envisions self-evolving systems that can discover new capabilities, transfer knowledge across domains, and operate within provable safety boundaries. Five frontiers define this phase. Cross-domain transfer evaluation: systems should demonstrate capability transfer across substantially different domains---from code to web, from simulation to physical robotics, from single-agent to multi-agent---with documented transfer efficiency, negative transfer rates, and preservation of core capabilities. Ecological archive sharing: organizations and research groups should share variant archives, evaluation trajectories, and failure reports through public infrastructure, enabling the community to build on collective experience rather than each starting from scratch. Provable safety boundaries: formal methods should provide guarantees for critical constraints---tool permissions, budget ceilings, data flow isolation, merge conditions, and rollback triggers---within a hybrid ``neural generation + symbolic guardrails + programmatic verification'' architecture.

Human-machine collaborative governance: humans and AI systems should collaboratively define risk levels, design evaluators, review edge cases, and calibrate safety boundaries, leveraging the complementary strengths of human judgment and machine-scale analysis. AGI engineering pathway: if the preceding frontiers are successfully navigated, the accumulated infrastructure---composable evolution loops, standardized evaluators, auditable memory, lineage tracking, safety sandboxes, and graduated autonomy---may constitute the engineering substrate needed for more general forms of artificial intelligence, where self-evolution is not a special-purpose mechanism but a fundamental capability of the system architecture.

Figure~\ref{fig:ch8-roadmap} visualizes this three-phase roadmap as a timeline.

% ====================================================================
% Figure 8.1: Research Roadmap (TikZ timeline)
% ====================================================================
\begin{figure}[htbp]
\centering
\begin{tikzpicture}[
  >=Stealth,
  phase/.style={
    rectangle, rounded corners=8pt, draw=#1!70!black, fill=#1!10,
    text width=4.2cm, minimum height=3.5cm, align=center, font=\small
  },
  arr/.style={-{Stealth[length=4mm]}, very thick, color=gray!50},
  note/.style={font=\scriptsize, color=gray!50!black, text width=3.8cm, align=left},
  timelabel/.style={font=\scriptsize\bfseries, color=gray!50},
]
% Phase 1: Foundation
\node[phase=cat1] (p1) at (0,0) {};
\node[font=\bfseries\small, color=cat1!80!black] at (0,1.3) {Phase 1};
\node[font=\bfseries\small, color=cat1!80!black] at (0,0.8) {Foundation};
\node[note] at (0,-0.1) {%
  $\bullet$ Standardized evaluator libraries\\
  $\bullet$ Auditable memory systems\\
  $\bullet$ Archive \& lineage infrastructure\\
  $\bullet$ Cost governance frameworks\\
  $\bullet$ Self-evolution reporting standards
};

% Phase 2: Graduated Autonomy
\node[phase=cat3] (p2) at (5.5,0) {};
\node[font=\bfseries\small, color=cat3!80!black] at (5.5,1.3) {Phase 2};
\node[font=\bfseries\small, color=cat3!80!black] at (5.5,0.8) {Graduated Autonomy};
\node[note] at (5.5,-0.1) {%
  $\bullet$ Low-risk autonomous improvement\\
  $\bullet$ Medium-risk gray-scale evaluation\\
  $\bullet$ High-risk human approval gates\\
  $\bullet$ Heterogeneous multi-agent checks\\
  $\bullet$ Auditable experience assetization
};

% Phase 3: Open-Ended Evolution
\node[phase=cat4] (p3) at (11,0) {};
\node[font=\bfseries\small, color=cat4!80!black] at (11,1.3) {Phase 3};
\node[font=\bfseries\small, color=cat4!80!black] at (11,0.8) {Open-Ended Evolution};
\node[note] at (11,-0.1) {%
  $\bullet$ Cross-domain transfer evaluation\\
  $\bullet$ Ecological archive sharing\\
  $\bullet$ Provable safety boundaries\\
  $\bullet$ Human-machine governance\\
  $\bullet$ AGI engineering pathway
};

% Arrows between phases
\draw[arr] (p1) -- (p2);
\draw[arr] (p2) -- (p3);

% Time labels
\node[timelabel] at (0,-2.4) {2026--2027};
\node[timelabel] at (5.5,-2.4) {2027--2029};
\node[timelabel] at (11,-2.4) {2029+};

% Timeline bar
\draw[very thick, gray!30] (-2.5,-2.1) -- (13.5,-2.1);
\foreach \x/\t in {0/Phase 1, 5.5/Phase 2, 11/Phase 3} {
  \draw[thick, gray!50] (\x,-2.25) -- (\x,-1.95);
}

\end{tikzpicture}
\caption{Phased roadmap for self-evolving AI. Phase 1 (2026--2027) builds foundation infrastructure: evaluators, memory, archives, and reporting standards. Phase 2 (2027--2029) introduces graduated autonomy and multi-agent checks-and-balances. Phase 3 (2029+) targets open-ended cross-domain evolution with provable safety boundaries.}
\label{fig:ch8-roadmap}
\end{figure}

Figure~\ref{fig:ch8-problems} presents a severity matrix of the most critical open problems identified in this survey, categorized by their impact on self-evolving systems and the urgency of addressing them.

% ====================================================================
% Figure 8.2: Open Problems Severity Matrix (TikZ)
% ====================================================================
\begin{figure}[htbp]
\centering
\begin{tikzpicture}
\def\cellwidth{3.0cm}
\def\cellheight{0.75cm}

% Title
\node[font=\small\bfseries] at (5.5, 4.2) {Open Problems Severity Matrix};

% Column headers
\node[font=\scriptsize\bfseries, text=gray!60] at (0, 3.6) {Problem Area};
\node[font=\scriptsize\bfseries, text=gray!60] at (3.5, 3.6) {Severity};
\node[font=\scriptsize\bfseries, text=gray!60] at (6.5, 3.6) {Urgency};
\node[font=\scriptsize\bfseries, text=gray!60] at (10, 3.6) {Current Status};

% Header underline
\draw[thick, gray!40] (-1.6, 3.3) -- (12.2, 3.3);

% Row 1: Evaluator Degeneration
\node[font=\scriptsize, anchor=west] at (-1.5, 2.8) {Evaluator Degeneration (Reward Hacking)};
\node[fill=cat5!60, minimum width=2cm, minimum height=\cellheight,
  rounded corners=2pt, font=\scriptsize\bfseries, text=white] at (3.5, 2.8) {Critical};
\node[fill=cat2!70, minimum width=2cm, minimum height=\cellheight,
  rounded corners=2pt, font=\scriptsize\bfseries, text=white] at (6.5, 2.8) {High};
\node[font=\scriptsize, anchor=west] at (8.5, 2.8) {Self-Rewarding limits emerging};

% Row 2: Self-Modification Safety
\node[font=\scriptsize, anchor=west] at (-1.5, 2.0) {Self-Modification Safety};
\node[fill=cat5!60, minimum width=2cm, minimum height=\cellheight,
  rounded corners=2pt, font=\scriptsize\bfseries, text=white] at (3.5, 2.0) {Critical};
\node[fill=cat5!60, minimum width=2cm, minimum height=\cellheight,
  rounded corners=2pt, font=\scriptsize\bfseries, text=white] at (6.5, 2.0) {Critical};
\node[font=\scriptsize, anchor=west] at (8.5, 2.0) {Docker/VM isolation insufficient};

% Row 3: Cross-Domain Transfer
\node[font=\scriptsize, anchor=west] at (-1.5, 1.2) {Cross-Domain Transfer};
\node[fill=cat2!70, minimum width=2cm, minimum height=\cellheight,
  rounded corners=2pt, font=\scriptsize\bfseries, text=white] at (3.5, 1.2) {High};
\node[fill=cat6!50, minimum width=2cm, minimum height=\cellheight,
  rounded corners=2pt, font=\scriptsize\bfseries, text=white] at (6.5, 1.2) {Medium};
\node[font=\scriptsize, anchor=west] at (8.5, 1.2) {Mostly benchmark-specific};

% Row 4: Multi-Agent Consensus Illusion
\node[font=\scriptsize, anchor=west] at (-1.5, 0.4) {Multi-Agent Consensus Illusion};
\node[fill=cat2!70, minimum width=2cm, minimum height=\cellheight,
  rounded corners=2pt, font=\scriptsize\bfseries, text=white] at (3.5, 0.4) {High};
\node[fill=cat2!70, minimum width=2cm, minimum height=\cellheight,
  rounded corners=2pt, font=\scriptsize\bfseries, text=white] at (6.5, 0.4) {High};
\node[font=\scriptsize, anchor=west] at (8.5, 0.4) {Debate cost high, convergence slow};

% Row 5: Evaluation Standardization
\node[font=\scriptsize, anchor=west] at (-1.5, -0.4) {Evaluation Standardization Gap};
\node[fill=cat2!70, minimum width=2cm, minimum height=\cellheight,
  rounded corners=2pt, font=\scriptsize\bfseries, text=white] at (3.5, -0.4) {High};
\node[fill=cat2!70, minimum width=2cm, minimum height=\cellheight,
  rounded corners=2pt, font=\scriptsize\bfseries, text=white] at (6.5, -0.4) {High};
\node[font=\scriptsize, anchor=west] at (8.5, -0.4) {Goodhart effect prevalent};

% Row 6: Cost Control
\node[font=\scriptsize, anchor=west] at (-1.5, -1.2) {Evolution Cost Control};
\node[fill=cat6!50, minimum width=2cm, minimum height=\cellheight,
  rounded corners=2pt, font=\scriptsize\bfseries, text=white] at (3.5, -1.2) {Medium};
\node[fill=cat2!70, minimum width=2cm, minimum height=\cellheight,
  rounded corners=2pt, font=\scriptsize\bfseries, text=white] at (6.5, -1.2) {High};
\node[font=\scriptsize, anchor=west] at (8.5, -1.2) {Must enter fitness function};

% Row 7: Negative Transfer / Forgetting
\node[font=\scriptsize, anchor=west] at (-1.5, -2.0) {Negative Transfer \& Forgetting};
\node[fill=cat6!50, minimum width=2cm, minimum height=\cellheight,
  rounded corners=2pt, font=\scriptsize\bfseries, text=white] at (3.5, -2.0) {Medium};
\node[fill=cat6!50, minimum width=2cm, minimum height=\cellheight,
  rounded corners=2pt, font=\scriptsize\bfseries, text=white] at (6.5, -2.0) {Medium};
\node[font=\scriptsize, anchor=west] at (8.5, -2.0) {Need regression test suites};

% Row 8: Memory Governance
\node[font=\scriptsize, anchor=west] at (-1.5, -2.8) {Memory Asset Governance};
\node[fill=cat6!50, minimum width=2cm, minimum height=\cellheight,
  rounded corners=2pt, font=\scriptsize\bfseries, text=white] at (3.5, -2.8) {Medium};
\node[fill=cat6!50, minimum width=2cm, minimum height=\cellheight,
  rounded corners=2pt, font=\scriptsize\bfseries, text=white] at (6.5, -2.8) {Medium};
\node[font=\scriptsize, anchor=west] at (8.5, -2.8) {Research stage};

% Bottom border
\draw[thick, gray!40] (-1.6, -3.3) -- (12.2, -3.3);

% Legend
\node[font=\tiny\bfseries, anchor=west] at (-1.5, -3.7) {Legend:};
\node[fill=cat5!60, minimum width=0.8cm, minimum height=0.35cm,
  rounded corners=1pt] at (1.0, -3.7) {};
\node[font=\tiny, anchor=west] at (1.5, -3.7) {Critical};
\node[fill=cat2!70, minimum width=0.8cm, minimum height=0.35cm,
  rounded corners=1pt] at (3.0, -3.7) {};
\node[font=\tiny, anchor=west] at (3.5, -3.7) {High};
\node[fill=cat6!50, minimum width=0.8cm, minimum height=0.35cm,
  rounded corners=1pt] at (4.8, -3.7) {};
\node[font=\tiny, anchor=west] at (5.3, -3.7) {Medium};

\end{tikzpicture}
\caption{Open problems severity matrix for self-evolving AI. Problems are classified by severity (impact on self-evolution systems) and urgency (how quickly they must be addressed). Color coding: red = critical, orange = high, cyan = medium.}
\label{fig:ch8-problems}
\end{figure}

\begin{table}[htbp]
\centering
\small
\begin{tabular}{p{27mm}p{34mm}p{40mm}p{38mm}}
\toprule
\textbf{Future Direction} & \textbf{Near-Term Actions} & \textbf{Key Metrics} & \textbf{Primary Risks} \\
\midrule
Automated Evaluators & Unit tests, environment replay, LLM judge calibration sets & Verifier coverage, false positive rate, drift rate, cost per verification & Evaluator overfitted or compromised by agent \\
Multi-Agent Collaboration & Generator/verifier/red-team separation, heterogeneous model routing & Independent error rate, dispute resolution rate, cost per success & Consensus illusion and redundant calls inflating cost \\
Safe Self-Evolution & Sandbox, permission minimization, lineage, rollback gating & High-risk operation interception rate, rollback time, privilege violation count & Permission boundaries bypassed via prompt injection or tool output poisoning \\
Cross-Domain Transfer & Experience libraries, task cluster routing, regression suites & New-domain cold-start samples, negative transfer rate, old-task retention & Benchmark-specific strategies fail to transfer \\
AGI Engineering Path & Open-ended archive, tool creation, governance auditing & New skill discovery rate, explainable lineage ratio, human audit burden & Capability growth outpaces verification and governance \\
\bottomrule
\end{tabular}
\caption{Engineering evaluation framework for future directions. Each direction requires concrete near-term actions, observable metrics, and risk gates to prevent research visions from degenerating into unverifiable narratives.}
\label{tab:ch8-direction-metrics}
\end{table}

% ====================================================================
% 8.7 Conclusion
% ====================================================================
\subsection{Conclusion}\label{sec:ch8-conclusion}

This survey has presented a comprehensive analysis of AI self-evolution, grounded in 107 core papers spanning 2022 to 2026, 348 GitHub repositories, 97 user pain points collected through Mom Test methodology, and a five-loop taxonomy that unifies the diverse landscape of self-improvement techniques. The central argument is that self-evolution is not a single technique but a family of composable mechanisms---specification-to-execution, search, evaluator, reflection, and population loops---each contributing distinct capabilities and each carrying specific limitations.

The field has made remarkable progress in demonstrating that self-improvement is possible. STaR showed that reasoning chains can be bootstrapped from self-generated solutions. Reflexion demonstrated that verbal reflection can substitute for parameter updates. ADAS proved that agent architectures can be automatically designed. DGM showed that open-ended self-modification can produce genuine capability gains. AlphaEvolve demonstrated that evolutionary search at Google scale can discover new algorithms and advance mathematical knowledge. Absolute Zero Reasoners proved that self-play with zero external data can bootstrap reasoning capabilities. AriadneMem and ReasoningBank showed that memory can be structured, versioned, and made auditable.

Yet the gap between research demonstration and production deployment remains wide. The Mom Test analysis revealed that the most common user complaints are not about capability deficits but about reliability, controllability, observability, and cost. Users do not want agents that are smarter in demos; they want agents that are reliable in production, that explain their decisions, that stay within budget, that can be rolled back when they fail, and that do not require constant human supervision to maintain their behavior. The star-analysis report reinforced this finding: projects with the highest community visibility often have the lowest production readiness scores, while the most rigorously engineered projects receive comparatively less attention.

The five-loop taxonomy provides a framework for understanding both the current landscape and the path forward. The specification-to-execution loop addresses the challenge of translating goals into executable pipelines. The search loop provides the exploration mechanism that discovers improvements. The evaluator loop supplies the feedback signal that guides search. The reflection loop converts failures into learning. The population loop maintains diversity and provides the stepping stones that enable open-ended discovery. No single system has successfully composed all five loops, and the interfaces between loops remain poorly standardized. Composability---the ability to mix and match loops based on task requirements, resource constraints, and safety profiles---is the key engineering challenge for the next phase of the field.

The safety and alignment challenge is particularly acute for self-evolving systems. Unlike static AI systems, self-evolving agents modify their own behavior, persist modifications in memory, and can amplify misaligned objectives across improvement cycles. The Misevolution analysis, the OEP experience poisoning attack, and the Capability Erosion study collectively demonstrate that safety risks are not hypothetical---they are concrete, documented, and potentially severe. Addressing these risks requires a layered approach: permission minimization, evolutionary sandboxing, auditable lineage, objective isolation, and graduated human governance. The aspiration for provable safety boundaries, while distant, provides a north star for engineering efforts.

Looking ahead, the maturation of self-evolving AI will be marked not by any single system claiming ``autonomous improvement'' but by the emergence of infrastructure that makes self-evolution safe, observable, reproducible, and composable. When an agent proposes a self-modification, we should be able to answer: Why was this modification proposed? What failure evidence motivated it? Which independent evaluators did it pass? What permissions and costs does it affect? How can it be rolled back if it fails? Does the improvement transfer to new tasks? When these questions have standard answers, supported by standard tools and standard protocols, self-evolving AI will have transitioned from a research curiosity to an engineering discipline. The roadmap proposed in this survey---foundation infrastructure, graduated autonomy, and open-ended evolution---provides a structured path toward that transition, grounded in the empirical evidence and community feedback that this survey has endeavored to synthesize.

For engineering teams, the practical starting point is small closed loops: select verifiable tasks, establish baselines, freeze safety boundaries, collect failure samples, build automated tests, and allow the agent to propose improvements only for low-risk components, with all improvements evaluated offline before gray-scale deployment. For research teams, the priority is reporting mechanism boundaries rather than only best-case scores: what fails, what regresses, what it costs, and where the system's competence ends. For the open-source community, the highest-value contributions are shared evaluators, failure samples, and reproducible evaluation trajectories rather than yet another agent framework. For enterprise governance teams, the core imperative is integrating agent self-evolution into existing change management, risk auditing, and cost control processes rather than treating it as a novel paradigm that operates outside established organizational structures. Only when these four constituencies---engineering, research, open-source, and governance---converge on shared standards, shared infrastructure, and shared accountability will self-evolving AI fulfill its promise as a reliable, controllable, and genuinely beneficial technology.
